% Provenance (never printed): see sources/outline_sources.md (Year 2: radiative forcing
% quantified; feedbacks and climate sensitivity; detection and attribution; scenarios; the
% remaining carbon budget; components of sea-level rise; physical science only). Every statement
% from IPCC AR6 WG1 (SPM; chapters 3, 5, 7 and 9; Annex III data), ledger rows in
% sources/ledger/book2.md (ipcc:*) and sources/ledger/book4.md (ar6*:*). Curves and numbers from
% figdata/bachelor-2/anthropogenic.py and the Book 4 second-writer scratch calculator calc/ch27.py.
\chapter{تغيّر المناخ الناجم عن الإنسان}\label{ch:b2:anthropogenic-climate-change}

منذ عام 1970 احترّ سطح الأرض أسرع من أيّ فترة أخرى من خمسين سنة على مدى ألفي سنة على
الأقلّ (\emph{بثقة عالية}). وقد روت مجلّدات المدرسة ما الذي تغيّر
ولماذا. ويحوّل هذا الفصل الرواية إلى ميزانية: كم واطًا لكلّ متر مربع تضيف المحرّكات
البشرية، وكم منها يعيده الكوكب الأدفأ إلى الفضاء، وكم من الاحترار يشتريه كلّ طنّ من ثنائي أكسيد
الكربون، وكيف يتقاسم ارتفاعَ البحر الماءُ الذي يحترّ والجليدُ الذي يذوب. ويأتي كلّ رقم
مع الشكّ فيه، باللغة المعايَرة للتقييمات.
% ledger: ipcc:fastest-2000yr, ipcc:likelihood

\begin{recall}
\textbf{من مجلّدات المدرسة}: التأثير الإشعاعي، و$\qty{3.93}{W.m^{-2}}$ لمضاعفة
ثنائي أكسيد الكربون، وقانونه اللوغاريتمي؛ والتغذية الراجعة المناخية؛ \omterm{def:b2:anthropogenic-climate-change:ecs}{وحساسية المناخ عند التوازن}، وأفضل
تقدير لها \qty{3}{\celsius}، و\emph{على نحو مرجح} من 2.5 إلى \qty{4}{\celsius}؛ \omterm{def:b2:anthropogenic-climate-change:da}{والإسناد}، وسيناريو الانبعاثات
والإسقاط؛ واللغة المعايَرة: \emph{مرجح} من 66 إلى \qty{100}{\percent}، و\emph{مرجح
جدًا} من 90 إلى \qty{100}{\percent}، والمدى $[x \text{ إلى } y]$ مدى مرجح جدًا. وقد
أدّى تأثير الإنسان إلى احترار المناخ بمقدار \qty{1.07}{\celsius} بين 1850--1900 و2010--2019.
\textbf{من كتب الفيزياء}: الجسم الأسود وقانون ستيفان--بولتزمان؛ والتمدّد الحراري.
\textbf{من كتب الأحياء}: دورة الكربون، والكسر المحمول جوًّا، ومصارف الكربون
وميزانية الكربون في الغلاف الحيوي. \textbf{من \cref{ch:b2:deep-time-climate}}: دورة الكربون البطيئة،
التي لا تستعيد فائضًا من ثنائي أكسيد الكربون إلا على مدى مئات الآلاف من السنين.
\end{recall}
% ledger: ipcc:erf-2xco2, ipcc:log-co2, ipcc:ecs, ipcc:likelihood, ipcc:attribution

\section{التأثير مقيسًا كمّيًا}\label{sec:b2:anthropogenic-climate-change:forcing}

\begin{definition}[التأثير الإشعاعي الفعّال]\label{def:b2:anthropogenic-climate-change:erf}
\emph{التأثير الإشعاعي الفعّال}\index{التأثير الإشعاعي الفعّال} (ERF) لاضطراب ما
--- تغيّر في أحد غازات الدفيئة، أو في الهباء الجوي، أو في الشمس --- هو التغيّر الذي يُحدثه في صافي
التدفّق الإشعاعي النازل عند أعلى الغلاف الجوي، بوحدة \unit{W.m^{-2}}، بعد أن يتكيّف معه الغلاف الجوي،
والسحب وسطح اليابسة، لكن قبل أن تستجيب درجة حرارة السطح العالمية. والتكيّفات المحسوبة هي
تلك التي يسبّبها العامل نفسه، بمعزل عن درجة حرارة
السطح.
\end{definition}
% ledger: ar6ch7:erf-def

يرسم التعريف الحدّ بين التأثير والاستجابة بما يتوقّف على درجة حرارة السطح، لا
بالسرعة. فالتأثير الآني لثنائي أكسيد الكربون هو تغيّر التدفّق بعد إضافة الغاز مباشرةً؛
والستراتوسفير، التي تبرد عندئذ في غضون أشهر، تغيّره مرّة أولى؛ ثم
تغيّره التروبوسفير وبخار الماء فيها وسحبها مرّة أخرى. ولمضاعفة ثنائي أكسيد
الكربون يجد التقييم \qty{3.75}{W.m^{-2}} بعد تكيّف الستراتوسفير،
وتكيّفات في التروبوسفير يلغي بعضها بعضًا جزئيًا: \qty{-0.60}{W.m^{-2}} من احترارها، و+0.22
من بخار الماء، و+0.45 من السحب و+0.11 من سطح اليابسة، أي \qty{+0.18}{W.m^{-2}} إجمالًا.
فيكون \omterm{def:b2:anthropogenic-climate-change:erf}{التأثير الإشعاعي الفعّال} $3.75 + 0.18 = \qty{3.93}{W.m^{-2}}$، أي أكثر بنسبة \qty{6}{\percent} من
القيمة \qty{3.71}{W.m^{-2}} في التقييم السابق (\cref{exo:b2:anthropogenic-climate-change:1}).
% ledger: ar6ch7:2xco2, ar6ch7:erf-def

أُعيد بناء تأثير كلّ عامل منذ عام 1750 سنةً بسنة. ويُظهره الشكل أدناه:
ثنائي أكسيد الكربون عند \qty{2.16}{W.m^{-2}} عام 2019، وغازات الدفيئة الأخرى --- الميثان،
وأكسيد النيتروز، والغازات المهلجنة والأوزون --- عند \qty{1.69}{W.m^{-2}} معًا،
والهباء الجوي عند \qty{-1.06}{W.m^{-2}}. والتأثير البشري الكلّي \qty{2.72}{W.m^{-2}}، ومداه
المرجح جدًا، من 1.96 إلى \qty{3.48}{W.m^{-2}}، واسع بسبب الهباء الجوي أساسًا. ويُقدَّر تأثيره الإشعاعي الفعّال على مدى
1750--2014 بالقيمة $-1.3$ [$-2.0$ إلى $-0.6$] \unit{W.m^{-2}} (\emph{بثقة متوسطة})،
ومعظمه عبر أثره في السحب، التي يزداد سطوعها حين يوفّر الهباء الجوي مزيدًا من القطيرات.
وتضيف الثورات البركانية قممًا سالبة قصيرة (\cref{ch:b2:eruptive-styles}): وأعمقها في
السجلّ، \qty{-4.8}{W.m^{-2}} عام 1816، يلي ثوران تامبورا عام 1815.
% ledger: data:b4-ar6-erf, ipcc:erf-total, ipcc:erf-agents, ar6ch7:aerosol, usgs:msh-comparisons

\begin{omfigure}
\centering
\begin{tikzpicture}
\begin{axis}[omrecord, width=0.92\linewidth, height=6.2cm, xmin=1750, xmax=2020, ymin=-5.2,
  ymax=3.8, xlabel={السنة}, ylabel={التأثير \babelsublr{ERF} منذ عام 1750 \babelsublr{(\unit{W.m^{-2}})}},
  xtick={1750,1800,...,2000}, /pgf/number format/1000 sep={}, xtick pos=bottom, ytick pos=left,
  tick label style={font=\scriptsize}, label style={font=\small},
  legend style={at={(0.5,-0.24)}, anchor=north, font=\scriptsize, draw=black!30,
  legend columns=3, /tikz/every even column/.append style={column sep=6pt}},
  legend cell align=left]
% computed in figdata/bachelor-2/anthropogenic.py part erf; ledger: data:b4-ar6-erf
\addplot[name path=hi, draw=none, forget plot] table[x=year, y=hi] {figdata/out/bachelor-2-anthropogenic-erf.dat};
\addplot[name path=lo, draw=none, forget plot] table[x=year, y=lo] {figdata/out/bachelor-2-anthropogenic-erf.dat};
\addplot[black!15, forget plot] fill between[of=hi and lo];
\addplot[black!45, line width=0.4pt] table[x=year, y=natural] {figdata/out/bachelor-2-anthropogenic-erf.dat};
\addlegendentry{\foreignlanguage{arabic}{طبيعي (البراكين، الشمس)}}
\addplot[omThm, line width=1pt] table[x=year, y=co2] {figdata/out/bachelor-2-anthropogenic-erf.dat};
\addlegendentry{ثنائي أكسيد الكربون}
\addplot[omThm!50!orange, line width=1pt, densely dashed] table[x=year, y=other_ghg] {figdata/out/bachelor-2-anthropogenic-erf.dat};
\addlegendentry{غازات الدفيئة الأخرى والأوزون}
\addplot[omDef, line width=1pt] table[x=year, y=aerosol] {figdata/out/bachelor-2-anthropogenic-erf.dat};
\addlegendentry{الهباء الجوي}
\addplot[black, line width=1.4pt] table[x=year, y=anthro] {figdata/out/bachelor-2-anthropogenic-erf.dat};
\addlegendentry{\foreignlanguage{arabic}{المجموع البشري (الرمادي: 5--95\,\%)}}
\draw[black!40] (axis cs:1750,0) -- (axis cs:2020,0);
\node[anchor=west, font=\scriptsize, text=black!60, inner sep=1pt] at (axis cs:1817.5,-4.75) {\foreignlanguage{arabic}{تامبورا $\leftarrow$}};
\end{axis}
\end{tikzpicture}
\par\smallskip{\small \omterm{def:b2:anthropogenic-climate-change:erf}{التأثير الإشعاعي الفعّال} من 1750 إلى 2019، نسبةً إلى 1750، بحسب المحرّك.
يبلغ التأثير البشري الكلّي \qty{2.72}{W.m^{-2}} عام 2019، مع تظليل المدى المرجح جدًا.
المعطيات: تقرير التقييم السادس للهيئة الحكومية الدولية المعنية بتغيّر المناخ، الفريق العامل الأول، المعطيات الموسّعة للملحق الثالث (\mbox{\textdir TLT CC~BY~4.0}).}
\end{omfigure}
% values: anthro 2.7204 (1.9569-3.4819), co2 2.1563, other_ghg 1.6850, aerosol -1.0583, volcanic -4.8162 in 1816 (anthropogenic.py numbers)

\section{ميزانية الطاقة}\label{sec:b2:anthropogenic-climate-change:energy}

يجعل التأثير الموجب الكوكب يكتسب طاقة. وكلّما احترّ أشعّ أكثر نحو الفضاء،
وتقلّص الكسب؛ ولا يستقرّ المناخ إلا حين يتوازن الاثنان.

\begin{definition}[اختلال توازن الطاقة الأرضية؛ معامل التغذية الراجعة المناخية]\label{def:b2:anthropogenic-climate-change:eei}
\emph{اختلال توازن الطاقة الأرضية}\index{اختلال توازن الطاقة الأرضية} $N$ هو صافي تدفّق الطاقة الداخل إلى
النظام المناخي عند أعلى الغلاف الجوي، متوسَّطًا على الكرة الأرضية وعلى سنة أو أكثر،
بوحدة \unit{W.m^{-2}}: ضوء الشمس الممتصّ ناقص الأشعّة تحت الحمراء الخارجة. وهو المعدّل الذي يكتسب به المحيط
واليابسة والجليد والهواء الحرارة، لكلّ متر مربع من سطح الأرض.
\end{definition}

\begin{definition}[معامل التغذية الراجعة المناخية]\label{def:b2:anthropogenic-climate-change:lambda}
\emph{معامل التغذية الراجعة المناخية}\index{معامل التغذية الراجعة المناخية} $\lambda$ هو التغيّر في
صافي التدفّق النازل عند أعلى الغلاف الجوي لكلّ درجة من احترار السطح العالمي، بوحدة
\unit{W.m^{-2}.\celsius^{-1}}، الناجم عن استجابة النظام المناخي لذلك الاحترار.
والمناخ المستقرّ له $\lambda < 0$: فالاحترار يزيد فقدان الطاقة نحو الفضاء.
\end{definition}

كتب مجلد المدرسة التوازن بمقدار $\lambda$، وهو عدد موجب؛ ومن هنا فصاعدًا
تُحفظ الإشارة، كما في التقييمات، التي تكتب هذا المعامل $\alpha$.
% ledger: ar6ch7:energy-budget

\begin{proposition}[توازن الطاقة العالمي]\label{prop:b2:anthropogenic-climate-change:balance}
ليكن $\Delta F$ \omterm{def:b2:anthropogenic-climate-change:erf}{التأثير الإشعاعي الفعّال} و$\Delta T$ احترار السطح العالمي
نسبةً إلى حالة توازن. إذا كانت استجابة التدفّق الخارج خطّية في $\Delta T$، فإن
\[
  \Delta N = \Delta F + \lambda\,\Delta T .
\]
وعند توازن جديد $\Delta N = 0$، فيكون $\Delta T_{\text{توازن}} = -\Delta F/\lambda$.
\end{proposition}
\begin{proof}
بحسب التعريف يغيّر التأثير صافي التدفّق النازل بمقدار $\Delta F$ عند درجة حرارة سطح
ثابتة، وتغيّره الاستجابة للاحترار بمقدار $\lambda\,\Delta T$ إلى الرتبة الأولى؛
والاختلال هو مجموعهما. وبجعله مساويًا للصفر نحصل على التوازن.
\end{proof}

قياس $N$ يعني قياس الحرارة التي يكتسبها النظام، والمحيط يحتفظ بها كلّها تقريبًا.
فمن 1971 إلى 2018 اكتسب النظام المناخي \qty{435}{ZJ} (زيتاجول، أي \qty{e21}{J}): المحيط
\qty{91}{\percent}، منها \qty{242}{ZJ} في الأمتار \qty{700}{m} العليا منه و\qty{154}{ZJ} تحتها؛ واليابسة
\qty{5}{\percent}، والجليد الذائب \qty{3}{\percent}، والغلاف الجوي \qty{1}{\percent}. وكان متوسط
الاختلال \qty{0.57}{W.m^{-2}} على مدى 1971--2018، وقد ازداد: \qty{0.79}{W.m^{-2}}
[0.52 إلى 1.06] على مدى 2006--2018 (\cref{exo:b2:anthropogenic-climate-change:2}).
% ledger: ar6ch7:inventory, ipcc:heating

\begin{omfigure}
\centering
\begin{tikzpicture}[font=\small, >=Stealth]
\draw[fill=omDef!8, draw=omDef!60] (-3.2,0) rectangle (3.2,1.6);
\node[align=center] at (0,1.05) {النظام المناخي};
\node[align=center, font=\scriptsize] at (0,0.5) {\foreignlanguage{arabic}{المحيط \qty{91}{\percent}، اليابسة \qty{5}{\percent}، الجليد \qty{3}{\percent}، الهواء \qty{1}{\percent}}};
\draw[->, line width=2.2pt, omThm] (-2.0,3.3) -- node[pos=0.12, left=3pt, align=right, font=\scriptsize] {\foreignlanguage{arabic}{التأثير $\Delta F$}} (-2.0,1.7);
\draw[->, line width=2.2pt, omDef] (2.0,1.7) -- node[pos=0.88, right=3pt, align=left, font=\scriptsize] {\foreignlanguage{arabic}{الاستجابة $\lambda\,\Delta T$}} (2.0,3.3);
\draw[->, line width=1.4pt, black!70] (0,-0.1) -- node[right=3pt, align=left, font=\scriptsize] {\foreignlanguage{arabic}{المخزَّن: $\Delta N = \Delta F + \lambda\,\Delta T$}} (0,-1.0);
\draw[black!40, dashed] (-3.4,2.4) -- (3.4,2.4) node[right, font=\scriptsize, text=black!60] {أعلى الغلاف الجوي};
\end{tikzpicture}
\par\smallskip{\small ميزانية الطاقة لمناخ خاضع لتأثير (تخطيطي). يضيف التأثير طاقة عند
أعلى الغلاف الجوي، وتنزع استجابة الكوكب الأدفأ جزءًا منها ($\lambda\,\Delta T$
سالبة)، ويُخزَّن الفرق، في المحيط أساسًا. حصص الحرارة المكتسبة من 1971
إلى 2018.}
\end{omfigure}
% ledger: ar6ch7:inventory

\begin{method}[حساسية مناخية من ميزانية الطاقة التاريخية]\label{met:b2:anthropogenic-climate-change:energy-budget}
\begin{enumerate}
  \item اختر فترتين، فترة مرجعية قريبة من التوازن وفترة حديثة؛ وخذ التغيّر في
        التأثير $\Delta F$، وفي الاختلال $\Delta N$ وفي درجة الحرارة العالمية $\Delta T$.
  \item حُلّ معادلة التوازن لإيجاد معامل التغذية الراجعة: $\lambda = (\Delta N - \Delta F)/\Delta T$.
  \item تكون الحساسية لمضاعفةٍ عندئذ $-F_{2\times}/\lambda$؛ والاحترار في وقت
        التأثير، مقيسًا إلى $F_{2\times}$، هو $F_{2\times}\Delta T/\Delta F$.
\end{enumerate}
\end{method}

من 1850--1900 إلى 2006--2019 يأخذ التقييم $\Delta N = \qty{0.59}{W.m^{-2}}$، و$\Delta T =
\qty{1.03}{\celsius}$ و$\Delta F = \qty{2.20}{W.m^{-2}}$. فيكون $\lambda = (0.59 - 2.20)/1.03 =
\qty{-1.56}{W.m^{-2}.\celsius^{-1}}$، والحساسية $3.93/1.56 = \qty{2.5}{\celsius}$،
وهي القيمة التي يجدها التقييم، بمدى مرجح جدًا من 1.6 إلى \qty{4.8}{\celsius}. لكن هذه
حساسية \emph{فعلية}، لتغذيات الاحترار الراجعة حتى الآن. فقد كان الاحترار
أشدّ حيث التغذيات الراجعة أشدّ سلبيةً، في غرب المحيط الهادئ المداري، وأضعف فوق
شرق المحيط الهادئ والمحيط الجنوبي؛ ومع تطوّر هذا النمط يُتوقّع أن يصبح معامل التغذية الراجعة
أقلّ سلبيةً (\emph{بثقة عالية})، فتكون الحساسية البعيدة المدى أعلى من
التقدير التاريخي (\cref{exo:b2:anthropogenic-climate-change:7}).
% ledger: ar6ch7:energy-budget, ar6ch7:pattern, ipcc:erf-2xco2
% values: lambda_energy -1.563, ecs_energy 2.514 (anthropogenic.py numbers)

\section{التغذيات الراجعة والحساسية}\label{sec:b2:anthropogenic-climate-change:feedbacks}

\begin{definition}[استجابة بلانك؛ التغذيات الراجعة لبخار الماء ولمعدّل الانخفاض الحراري وللسحب]\label{def:b2:anthropogenic-climate-change:feedbacks}
يُجزَّأ معامل التغذية الراجعة إلى إسهامات الكمّيات التي تتغيّر مع درجة حرارة
السطح. \emph{استجابة بلانك}\index{استجابة بلانك} هي الأشعّة تحت الحمراء الإضافية
المنبعثة نحو الفضاء حين يحترّ السطح والتروبوسفير كلّها بالمقدار نفسه.
و\emph{التغذية الراجعة لبخار الماء}\index{التغذية الراجعة لبخار الماء} هي تغيّر التدفّق الناجم عن ازدياد
بخار الماء الجوّي مع درجة الحرارة.
و\emph{التغذية الراجعة لمعدّل الانخفاض الحراري}\index{التغذية الراجعة لمعدّل الانخفاض الحراري} هي تغيّر التدفّق الناجم عن احترار غير منتظم مع الارتفاع. و\emph{التغذية
الراجعة للسحب}\index{التغذية الراجعة للسحب} هي تغيّر التدفّق الناجم عن تغيّر كمّية السحب وارتفاعها
وسطوعها.
\end{definition}

\begin{proposition}[استجابة بلانك]\label{prop:b2:anthropogenic-climate-change:planck}
إذا كانت الأرض تبعث نحو الفضاء كسرًا ثابتًا $1 - g$ من الأشعّة تحت الحمراء $\sigma T^4$ لسطحها عند
درجة الحرارة $T$، فإن \omterm{def:b2:anthropogenic-climate-change:feedbacks}{استجابة بلانك} لها هي
\[
  \lambda_{\text P} = -(1 - g)\,4\sigma T^3 .
\]
ومع $g = 0.4$ و$T = \qty{288}{K}$، $\lambda_{\text P} = \qty{-3.25}{W.m^{-2}.K^{-1}}$.
\end{proposition}
\begin{proof}
التدفّق الخارج هو $(1 - g)\sigma T^4$؛ ومشتقّه بالنسبة إلى $T$ هو $(1 - g)\,4\sigma
T^3$، والفقد الإضافي إسهام سالب في صافي التدفّق النازل. وعدديًا $0.6
\times 4 \times \num{5.670e-8} \times 288^3 = 3.25$.
\end{proof}
% ledger: ar6ch7:planck, const:sigma, ipcc:budget-surface
% values: planck -3.2509 (anthropogenic.py numbers)

الكسر $g$ هو أثر الدفيئة المعيَّر: فالغلاف الجوي يحتفظ بالقيمة \qty{159}{W.m^{-2}} من
\qty{398}{W.m^{-2}} التي يبعثها السطح، $g \approx 0.4$. ويجد التقييم، بشيفرات إشعاعية
مفصّلة، $-3.22$ [$-3.4$ إلى $-3.0$] (\emph{بثقة عالية}). ويتبع بخار الماء
قانون كلاوزيوس--كلابيرون على المقياس العالمي، ومزيد من البخار يمتصّ مزيدًا من الأشعّة تحت الحمراء: \omterm{def:b2:anthropogenic-climate-change:feedbacks}{فالتغذية
الراجعة لبخار الماء} موجبة بقوّة، \qty{1.85}{W.m^{-2}.\celsius^{-1}} من الأقمار الصناعية.
وفي المنطقة المدارية تحترّ التروبوسفير العليا أكثر من السطح، فتشعّ أكثر: \omterm{def:b2:anthropogenic-climate-change:feedbacks}{فالتغذية
الراجعة لمعدّل الانخفاض الحراري} سالبة هناك، وموجبة قرب القطبين، ونحو $-0.50$ عالميًا. ويُقيَّم الاثنان
معًا، \num{1.30} [1.1 إلى 1.5]، لأن ما يضعه نموذج في أحدهما يميل إلى
أخذه من الآخر. ويُعتم الثلج والجليد الذائبان السطح، +0.35؛ وتضيف السحب \num{0.42}
[$-0.10$ إلى $0.94$]، وهي أكبر مصادر الشكّ. وأن تكون التغذية الراجعة الصافية للسحب سالبةً أمرٌ \emph{غير مرجح جدًا}.
% ledger: ar6ch7:planck, ar6ch7:wv-lr, ar6ch7:table710, ipcc:cloud-negative, ipcc:budget-surface

\begin{omfigure}
\centering
\begin{tikzpicture}
\begin{axis}[xbar, width=0.78\linewidth, height=5.4cm, xmin=-3.7, xmax=1.8, ymin=0.4, ymax=6.6,
  ytick={1,...,6}, yticklabels={الصافي, الجيوفيزيائية الحيوية والكيميائية, السحب, ألبيدو السطح,
  بخار الماء + معدّل الانخفاض, بلانك}, y tick label style={font=\scriptsize},
  tick label style={font=\scriptsize}, xlabel={معامل التغذية الراجعة \babelsublr{(\unit{W.m^{-2}.\celsius^{-1}})}},
  label style={font=\small}, bar width=7pt, axis x line=bottom, axis y line=left,
  axis line style={-}, xmajorgrids, grid style={gray!22, line width=0.3pt}, extra x ticks={0},
  extra x tick style={grid style={black!60, line width=0.5pt}}]
% ledger: ar6ch7:table710 (central values and very likely ranges)
\addplot[fill=omDef!55, draw=omDef!80!black, error bars/.cd, x dir=both, x explicit,
  error bar style={black!70}] coordinates {
  (-3.22,6) += (0.22,0) -= (0.18,0)
  (-0.01,2) += (0.26,0) -= (0.26,0)
  (-1.16,1) += (0.65,0) -= (0.65,0)};
\addplot[fill=omThm!55, draw=omThm!80!black, error bars/.cd, x dir=both, x explicit,
  error bar style={black!70}] coordinates {
  (1.30,5) += (0.20,0) -= (0.20,0)
  (0.35,4) += (0.25,0) -= (0.25,0)
  (0.42,3) += (0.52,0) -= (0.52,0)};
\end{axis}
\end{tikzpicture}
\par\smallskip{\small التغذيات الراجعة المناخية كما قيّمتها الهيئة الحكومية الدولية المعنية بتغيّر المناخ (تقرير التقييم السادس، الفريق العامل
الأول، الجدول 7.10)، القيم المركزية والمديات المرجحة جدًا. التغذيات الراجعة الموجبة تضخّم الاحترار
(أحمر)، والسالبة تخمده (أزرق)؛ والصافي هو المجموع.}
\end{omfigure}

\begin{definition}[حساسية المناخ عند التوازن]\label{def:b2:anthropogenic-climate-change:ecs}
\emph{حساسية المناخ عند التوازن}\index{حساسية المناخ عند التوازن} (ECS) هي احترار السطح
العالمي، عند التوازن الجديد، الناجم عن مضاعفة ثنائي أكسيد الكربون في
الهواء انطلاقًا من مستواه قبل الصناعة، مع إبقاء الأغطية الجليدية ثابتة.
\end{definition}

\begin{proposition}[الحساسية والتغذية الراجعة]\label{prop:b2:anthropogenic-climate-change:ecs}
إذا كان معامل التغذية الراجعة ثابتًا، فإن $\text{ECS} = -F_{2\times}/\lambda$. ومع $F_{2\times} =
\qty{3.93}{W.m^{-2}}$ والقيمة الصافية $\lambda = \qty{-1.16}{W.m^{-2}.\celsius^{-1}}$ للتغذيات الراجعة
أعلاه، $\text{ECS} = \qty{3.4}{\celsius}$؛ أمّا \omterm{def:b2:anthropogenic-climate-change:feedbacks}{استجابة بلانك} وحدها فكانت ستعطي $3.93/3.22 =
\qty{1.2}{\celsius}$.
\end{proposition}
\begin{proof}
طبّق \cref{prop:b2:anthropogenic-climate-change:balance} عند التوازن مع $\Delta F =
F_{2\times}$.
\end{proof}
% ledger: ar6ch7:table710, ipcc:erf-2xco2
% values: ecs_process 3.388 (anthropogenic.py numbers)

التغذيات الراجعة خطّ واحد من خطوط الأدلّة؛ والاحترار منذ 1850، وميزانية الطاقة ومناخات
الماضي خطوط أخرى. ومجتمعةً، تعطي \omterm{def:b2:anthropogenic-climate-change:ecs}{حساسية المناخ عند التوازن} المقيّمة \qty{3}{\celsius}، و\emph{على نحو مرجح}
من 2.5 إلى \qty{4}{\celsius} (\emph{بثقة عالية})، و\emph{على نحو مرجح جدًا} من 2 إلى \qty{5}{\celsius}.
ولأن هذه الحساسية تتغيّر بحسب $1/\lambda$، فإن شكًّا متناظرًا في $\lambda$ يعطي مدى
للحساسية ممطوطًا نحو القيم المرتفعة (\cref{exo:b2:anthropogenic-climate-change:12}).
% ledger: ipcc:ecs

\begin{definition}[الاستجابة المناخية العابرة]\label{def:b2:anthropogenic-climate-change:tcr}
\emph{الاستجابة المناخية العابرة}\index{الاستجابة المناخية العابرة} (TCR) هي احترار السطح
العالمي في اللحظة التي يتضاعف فيها ثنائي أكسيد الكربون في الهواء، حين يرتفع بنسبة \qty{1}{\percent}
في السنة انطلاقًا من مستواه قبل الصناعة، أي بعد نحو سبعين سنة من البداية.
\end{definition}

الاستجابة العابرة أصغر من حساسية التوازن لأن المحيط يظلّ يمتصّ الحرارة حين تُبلَغ
المضاعفة. وأفضل تقدير لها \qty{1.8}{\celsius}، و\emph{على نحو مرجح} من 1.4 إلى \qty{2.2}{\celsius}
(\emph{بثقة عالية})؛ وميزانية الطاقة التاريخية وحدها تعطي $3.93 \times 1.03/2.20 =
\qty{1.8}{\celsius}$ أيضًا.
% ledger: ar6ch7:tcr, ar6ch7:energy-budget
% values: tcr_energy 1.840 (anthropogenic.py numbers)

\section{الكشف والإسناد}\label{sec:b2:anthropogenic-climate-change:attribution}

\begin{definition}[الكشف؛ الإسناد؛ البصمة المناخية]\label{def:b2:anthropogenic-climate-change:da}
\emph{الكشف}\index{الكشف (عن تغيّر المناخ)} عن تغيّر في المناخ هو بيان أن
المناخ قد تغيّر بمعنى إحصائي محدّد، يتجاوز ما يمكن أن تُحدثه تقلّبيته الداخلية،
من دون تحديد السبب. و\emph{الإسناد}\index{إسناد} هو تقييم
إسهامات عدّة أسباب في تغيّر مكتشَف، مع احتمال محدّد. و\emph{البصمة
المناخية}\index{البصمة المناخية} هي النمط المكاني والزماني للاستجابة لمحرّك
واحد، كما تحاكيه نماذج المناخ، والذي يُتعرَّف به على ذلك المحرّك في الأرصاد.
\end{definition}

المحرّكات المختلفة تترك بصمات مختلفة. فمزيد من غازات الدفيئة يُدفئ التروبوسفير ويبرّد
الستراتوسفير، لأن غطاءً أسمك من الأشعّة تحت الحمراء يجعل الطبقات العليا تشعّ أكثر نحو الفضاء؛
أمّا شمس أشدّ سطوعًا فكانت ستُدفئ الاثنين. والأرصاد تُظهر النمط الأوّل: فقد احترّت التروبوسفير
منذ خمسينيات القرن العشرين على الأقلّ، ومن \emph{المؤكد تقريبًا} أن الستراتوسفير قد بردت. ومن
\emph{المرجح جدًا} أن غازات الدفيئة البشرية كانت المحرّك الرئيسي لاحترار التروبوسفير
منذ 1979، ومن \emph{المرجح للغاية} أن استنزاف طبقة الأوزون كان المحرّك الرئيسي
لتبرّد الستراتوسفير السفلى من 1979 إلى منتصف تسعينيات القرن العشرين.
% ledger: ar6ch3:upper-air

\begin{method}[البصمة المثلى]\label{met:b2:anthropogenic-climate-change:fingerprinting}
\begin{enumerate}
  \item حاكِ الاستجابة لكلّ محرّك على حدة: البصمات $x_1, \dots, x_k$، على هيئة
        متّجهات من القيم في المكان والزمان.
  \item اكتب الأرصاد على الصورة $y = \sum_i \beta_i x_i + \varepsilon$، حيث $\varepsilon$
        التقلّبية الداخلية، التي تؤخذ إحصاءاتها من محاكاة طويلة غير خاضعة لتأثير.
  \item قدّر معاملات القياس $\beta_i$ بالمربّعات الصغرى، موزونةً بمقلوب
        التقلّبية بحيث تكون للمركّبات الأكثر ضجيجًا أقلّ وزن، مع مجالات الثقة الخاصّة بها.
  \item قيمة $\beta_i$ أكبر من الصفر بدلالة إحصائية تعني أن الاستجابة للمحرّك $i$
        مكتشَفة؛ ومجال يحتوي على 1 ويستبعد 0 يعني أن الاستجابة المحاكاة
        متوافقة مع الأرصاد، وأن التغيّر يمكن إسناده جزئيًا إلى ذلك المحرّك.
\end{enumerate}
\end{method}
% ledger: ar6ch3:fingerprint

إذا طُبّقت الطريقة على درجة حرارة السطح، أعطت \omterm{def:b2:anthropogenic-climate-change:da}{الإسناد} المذكور في البداية: فمن
\qty{1.06}{\celsius} المرصودة بين 1850--1900 و2010--2019، سبّب تأثير الإنسان \emph{على نحو مرجح}
من 0.8 إلى \qty{1.3}{\celsius}، وأفضل تقدير \qty{1.07}{\celsius}؛ ولا تُسهم المحرّكات الطبيعية
والتقلّبية الداخلية بأكثر من بضعة أعشار الدرجة في أيّ من الاتجاهين
(\cref{exo:b2:anthropogenic-climate-change:9,exo:b2:anthropogenic-climate-change:10}).
% ledger: ipcc:attribution, ipcc:observed-2010-2019

\section{السيناريوهات والانبعاثات التراكمية}\label{sec:b2:anthropogenic-climate-change:scenarios}

\begin{definition}[المسار الاجتماعي الاقتصادي المشترك]\label{def:b2:anthropogenic-climate-change:ssp}
\emph{المسار الاجتماعي الاقتصادي المشترك}\index{المسار الاجتماعي الاقتصادي المشترك} (SSP) هو واحد من خمس
روايات لكيفية تطوّر السكّان والاقتصادات والتكنولوجيا خلال القرن.
وسيناريوهات الانبعاثات في التقييم تُسمّى SSP$x$-$y$: حيث يسمّي $x$ المسار، و$y$
التأثير الإشعاعي التقريبي، بوحدة \unit{W.m^{-2}}، الذي يبلغه السيناريو عام 2100.
\end{definition}
% ledger: ar6spm:ssp

تمتدّ السيناريوهات الخمسة الرئيسية من SSP1-1.9 إلى SSP5-8.5. ويأخذها تقييم العلوم الفيزيائية
كما هي: فهو لا يزن احتمالها ولا يقيّم الروايات. ولكلّ منها، تعطي نماذج المناخ
احترارًا مُسقَطًا للفترة 2081--2100 فوق 1850--1900، من \qty{1.4}{\celsius} (على نحو مرجح جدًا من 1.0
إلى 1.8) في السيناريو SSP1-1.9 إلى \qty{4.4}{\celsius} (من 3.3 إلى 5.7) في السيناريو SSP5-8.5. وتشترك الإسقاطات في
قانون بسيط.
% ledger: ar6spm:ssp, ipcc:proj-temp

\begin{definition}[الاستجابة المناخية العابرة للانبعاثات التراكمية من CO$_2$]\label{def:b2:anthropogenic-climate-change:tcre}
\emph{الاستجابة المناخية العابرة للانبعاثات التراكمية من \ce{CO2}}\index{الاستجابة المناخية العابرة للانبعاثات التراكمية من CO$_2$}
(TCRE) هي احترار السطح العالمي لكلّ وحدة من ثنائي أكسيد
الكربون المنبعث إجمالًا منذ فترة ما قبل الصناعة، بوحدة \unit{\celsius} لكلّ
\qty{1000}{Gt} من \ce{CO2}.
\end{definition}

\begin{proposition}[الاحترار يتبع الانبعاثات التراكمية]\label{prop:b2:anthropogenic-climate-change:tcre}
بتقريب جيّد يتناسب الاحترار الذي يسبّبه ثنائي أكسيد الكربون مع إجمالي ما انبعث منه،
أيًّا كان مسار الانبعاثات: $\Delta T \approx \text{TCRE} \times E$، حيث $E$ الانبعاثات
التراكمية.
\end{proposition}
\admitted
% ledger: ar6spm:tcre, ar6ch5:tcre-linear

هذا القانون ليس بديهيًا. فالتأثير لا يزداد إلا كلوغاريتم التركيز، فكلّ
طنّ مضاف يُدفئ أقلّ؛ لكن المحيط يمتصّ الحرارة والكربون بكفاءة تتناقص باستمرار مع
تراكم الانبعاثات، فيبقى مزيد من كلّ طنّ في الهواء ويتأخّر احتراره أقلّ.
والأثران يلغي أحدهما الآخر تقريبًا، وهذا ما تتّفق جميع الدراسات على أنه أصل القانون. والعلاقة
شبه خطّية (\emph{بثقة عالية})، وكلّ \qty{1000}{Gt} من ثنائي أكسيد الكربون تسبّب \emph{على نحو مرجح}
من 0.27 إلى \qty{0.63}{\celsius} من الاحترار، وأفضل تقدير \qty{0.45}{\celsius}، أي \qty{1.65}{\celsius}
لكلّ \qty{1000}{PgC} من الكربون (\cref{exo:b2:anthropogenic-climate-change:4}). ومن 1850 إلى 2019،
انبعث $\num{2390} \pm \qty{240}{Gt}$ من ثنائي أكسيد الكربون.
% ledger: ar6ch5:tcre-linear, ar6spm:tcre, ar6spm:budget

\begin{definition}[ميزانية الكربون المتبقّية]\label{def:b2:anthropogenic-climate-change:budget}
\emph{ميزانية الكربون المتبقّية}\index{ميزانية الكربون المتبقّية} لحدّ معيّن من درجة الحرارة هي
كمّية ثنائي أكسيد الكربون التراكمية التي لا يزال ممكنًا إطلاقها، ابتداءً من تاريخ محدّد، مع بقاء الاحترار العالمي
دون ذلك الحدّ باحتمال محدّد، مع أخذ الاحترار الذي تسبّبه المحرّكات البشرية الأخرى
في الحسبان.
\end{definition}

\begin{method}[ميزانية الكربون المتبقّية انطلاقًا من الاستجابة TCRE]\label{met:b2:anthropogenic-climate-change:budget}
\begin{enumerate}
  \item اطرح من الحدّ الاحترار البشري الذي بُلغ حتى الآن، \qty{1.07}{\celsius}
        للفترة 2010--2019.
  \item اطرح الاحترار الذي ستضيفه المحرّكات الأخرى حين تبلغ انبعاثات ثنائي أكسيد الكربون
        صافي الصفر.
  \item اقسم على الاستجابة TCRE: فهذه هي الميزانية ابتداءً من منتصف 2010--2019.
  \item اطرح الانبعاثات منذ ذلك الحين، وهامشًا للتغذيات الراجعة البطيئة لنظام الأرض التي
        تُغفلها النماذج (ذوبان التربة الصقيعية أكبرها).
\end{enumerate}
\end{method}

يجمع التقييم خمس مركّبات على هذا النحو: الاستجابة TCRE؛ والاحترار التاريخي؛
واحترارًا إضافيًا من المحرّكات الأخرى قدره نحو 0.1 إلى \qty{0.2}{\celsius}؛ والاحترار الذي قد
يستمرّ بعد توقّف الانبعاثات، المقيَّم بصفر في المتوسط (\emph{على نحو مرجح} في حدود $\pm
\qty{0.19}{\celsius}$)؛ والتغذيات الراجعة لنظام الأرض، نحو \qty{26}{Gt} من ثنائي أكسيد الكربون لكلّ
درجة. ولحدّ \qty{1.5}{\celsius}، أي \qty{0.43}{\celsius} فوق 2010--2019، تبلغ الميزانية المتبقّية ابتداءً من
مطلع 2020 \qty{500}{Gt} من ثنائي أكسيد الكربون باحتمال \qty{50}{\percent}،
و\qty{1350}{Gt} لحدّ \qty{2}{\celsius}. والخيارات المتّخذة بشأن المحرّكات الأخرى تحرّك هذه القيم بمقدار
\qty{220}{Gt} أو أكثر؛ وتعيد مسألة نهاية الأسبوع بناء القيمة الأولى.
% ledger: ar6ch5:budget-parts, ar6spm:budget

\begin{dated}{2025}
% ledger: data:co2-mlo, data:gistemp (values read from the datasets as fetched 2026-10-22; offset computed)
بلغ المتوسط السنوي لثنائي أكسيد الكربون في ماونا لوا \qty{427.35}{ppm} عام 2025، أي تأثيرًا قدره
$3.93\log_2(427.35/278) = \qty{2.44}{W.m^{-2}}$ بالقانون اللوغاريتمي. وفي سجلّ درجات الحرارة
لدى ناسا كانت 2024 أدفأ سنة، أعلى بمقدار \qty{1.50}{\celsius} من متوسط 1880--1900، وهي أقدم
سنوات ذلك السجلّ؛ وكانت 2025 عند \qty{1.40}{\celsius}.
\end{dated}

\begin{omfigure}
\centering
\begin{tikzpicture}
\begin{axis}[omrecord, width=0.86\linewidth, height=5.8cm, xmin=0, xmax=4000, ymin=0, ymax=2.6,
  xlabel={الانبعاثات التراكمية من \ce{CO2} منذ 1850 \babelsublr{(\unit{Gt})}},
  ylabel={الاحترار منذ 1850--1900 \babelsublr{(\unit{\celsius})}}, xtick={0,500,...,4000},
  /pgf/number format/1000 sep={\,}, xtick pos=bottom, ytick pos=left, clip=false,
  tick label style={font=\scriptsize}, label style={font=\small}]
% computed in figdata/bachelor-2/anthropogenic.py part tcre; ledger: ar6spm:tcre, ar6spm:budget
\addplot[name path=hi, draw=none, forget plot] table[x=E_GtCO2, y=dT_hi] {figdata/out/bachelor-2-anthropogenic-tcre.dat};
\addplot[name path=lo, draw=none, forget plot] table[x=E_GtCO2, y=dT_lo] {figdata/out/bachelor-2-anthropogenic-tcre.dat};
\addplot[omThm!18, forget plot] fill between[of=hi and lo];
\addplot[omThm, line width=1.3pt] table[x=E_GtCO2, y=dT_mid] {figdata/out/bachelor-2-anthropogenic-tcre.dat};
\draw[densely dashed, black!60] (axis cs:0,1.5) -- (axis cs:4000,1.5);
\draw[densely dashed, black!60] (axis cs:0,2.0) -- (axis cs:4000,2.0);
\node[anchor=north west, font=\scriptsize, text=black!70] at (axis cs:0,1.5) {\qty{1.5}{\celsius}};
\node[anchor=north west, font=\scriptsize, text=black!70] at (axis cs:0,2.0) {\qty{2}{\celsius}};
\fill[black] (axis cs:2390,1.07) circle (2.2pt);
\draw[black, line width=0.8pt] (axis cs:2150,1.07) -- (axis cs:2630,1.07);
\draw[black, line width=0.8pt] (axis cs:2390,0.8) -- (axis cs:2390,1.3);
\node[anchor=north west, font=\scriptsize, align=left] at (axis cs:2440,1.02) {\foreignlanguage{arabic}{الفترة 1850--2019:}\\\foreignlanguage{arabic}{الانبعاثات \qty{2390}{Gt}، والاحترار \qty{1.07}{\celsius}}};
\draw[<->, omDef, thick] (axis cs:2390,1.62) -- (axis cs:2890,1.62);
\node[anchor=south, font=\scriptsize, text=omDef, align=center] at (axis cs:2640,1.64) {\foreignlanguage{arabic}{تقدير \babelsublr{AR6}: \qty{500}{Gt}}};
\node[anchor=north west, font=\scriptsize, text=omThm, align=left] at (axis cs:150,2.5) {\foreignlanguage{arabic}{المدى المرجح للاستجابة TCRE،}\\\foreignlanguage{arabic}{من 0.27 إلى \qty{0.63}{\celsius} لكلّ \qty{1000}{Gt}}};
\end{axis}
\end{tikzpicture}
\par\smallskip{\small الاحترار مقابل الانبعاثات التراكمية لثنائي أكسيد الكربون للاستجابة TCRE المقيّمة،
\qty{0.45}{\celsius} لكلّ \qty{1000}{Gt} (الخطّ) ومداها المرجح (المظلّل)، مع الاحترار البشري
والانبعاثات حتى اليوم (النقطة، والمديات المرجحة على هيئة أشرطة). ويشير السهم إلى الميزانية المتبقّية لحدّ
\qty{1.5}{\celsius} باحتمال \qty{50}{\percent} ابتداءً من 2020: وهي أصغر ممّا تعطيه قراءة الخطّ وحده،
لأن المحرّكات الأخرى تُدفئ أيضًا.}
\end{omfigure}


\section{مستوى سطح البحر}\label{sec:b2:anthropogenic-climate-change:sea-level}

\begin{definition}[ميزانية مستوى سطح البحر؛ الارتفاع الحراري لمستوى سطح البحر]\label{def:b2:anthropogenic-climate-change:sea-level}
\emph{ميزانية مستوى سطح البحر}\index{ميزانية مستوى سطح البحر} هي مقارنة التغيّر المرصود في
المتوسط العالمي لمستوى سطح البحر بمجموع إسهاماته المقدّرة: تمدّد المحيط
الذي يحترّ، وذوبان الأنهار الجليدية والغطاءين الجليديين لغرينلاند والقارّة القطبية الجنوبية، وتغيّر
الماء المخزَّن على اليابسة. و\emph{الارتفاع الحراري لمستوى سطح البحر}\index{الارتفاع الحراري لمستوى سطح البحر} هو
الجزء الناجم عن التمدّد الحراري لماء البحر، عند كتلة ثابتة.
\end{definition}

\begin{proposition}[التمدّد الحراري لطبقة]\label{prop:b2:anthropogenic-climate-change:thermosteric}
طبقة من المحيط سمكها $H$، ومعامل تمدّدها الحراري $\alpha$، تحترّ بانتظام
بمقدار $\Delta T$ عند كتلة ومساحة ثابتتين، ترتفع بمقدار $\Delta h = \alpha\,\Delta T\,H$.
\end{proposition}
\begin{proof}
الكتلة $\rho H A$ لعمود مساحته $A$ ثابتة، فيكون $\Delta H/H = -\Delta\rho/\rho =
\alpha\,\Delta T$ إلى الرتبة الأولى، مع $\alpha = -(1/\rho)\,\partial\rho/\partial T$.
\end{proof}
% ledger: teos10:alpha-beta

في المحيط الحقيقي يزداد $\alpha$ مع درجة الحرارة والضغط، ويتركّز الاحترار قرب
السطح. ويلخّص التقييم النتيجة في كفاءة تمدّد قدرها \qty{0.113}{m} لكلّ
\qty{e24}{J} من الحرارة المضافة إلى المحيط، أو \qty{0.617}{m} لكلّ درجة من احترار المحيط كلّه
(\emph{بثقة متوسطة}). وعلى متوسط عمقه البالغ \qty{3682}{m}، يعطي ذلك قيمة فعلية للمعامل $\alpha$ قدرها
$0.617/3682 = \qty{1.68e-4}{\celsius^{-1}}$ (\cref{exo:b2:anthropogenic-climate-change:8}).
% ledger: ar6ch9:expansion, charette:ocean-volume

تُغلَق الميزانية، المرسومة في الشكل أدناه، ضمن حدود الشكّ فيها. فعلى مدى 1971--2018 يبلغ مجموع الإسهامات
\qty{2.00}{mm.yr^{-1}}، مقابل ارتفاع مرصود قدره \qty{2.33}{mm.yr^{-1}}؛ وعلى مدى 2006--2018
يبلغ مجموعها 3.61 مقابل 3.69. وكان التمدّد الحراري أكبر حدّ منفرد، نصف المجموع على مدى 1971--2018، لكن الجليد
يلحق به: فعلى مدى 2006--2018 أعطت الأنهار الجليدية والأغطية الجليدية معًا \qty{1.62}{mm.yr^{-1}}، أي أكثر من
التمدّد. ويظلّ المحيط أيضًا يحترّ في العمق زمنًا طويلًا بعد أن يستقرّ السطح، والأغطية
الجليدية تستجيب على مدى قرون؛ فمستوى سطح البحر ملتزم بالارتفاع لقرون إلى آلاف السنين وسيبقى
مرتفعًا لآلاف السنين (\emph{بثقة عالية}). وعلى مدى الألفي سنة القادمة
يُتوقّع أن يرتفع بنحو 2 إلى \qty{3}{m} إذا بقي الاحترار عند \qty{1.5}{\celsius}، ومن 2 إلى
\qty{6}{m} عند \qty{2}{\celsius} ومن 19 إلى \qty{22}{m} عند \qty{5}{\celsius} (\emph{بثقة منخفضة}).
% ledger: ar6ch9:table95, ar6spm:sea-commit

\begin{omfigure}
\centering
\begin{tikzpicture}
\begin{axis}[ybar stacked, width=0.72\linewidth, height=6.2cm, ymin=0, ymax=4.3, bar width=26pt,
  xtick=data, symbolic x coords={1971--2018, 1993--2018, 2006--2018}, enlarge x limits=0.25,
  ylabel={معدّل الارتفاع \babelsublr{(\unit{mm.yr^{-1}})}}, tick label style={font=\scriptsize},
  label style={font=\small}, ytick pos=left, xtick pos=bottom,
  legend style={at={(1.03,0.5)}, anchor=west, font=\scriptsize, draw=black!30},
  legend cell align=left, reverse legend]
% ledger: ar6ch9:table95
\addplot[fill=omDef!70, draw=omDef!90!black] coordinates {(1971--2018,1.01) (1993--2018,1.31) (2006--2018,1.39)};
\addlegendentry{التمدّد الحراري}
\addplot[fill=omDef!30, draw=omDef!70!black] coordinates {(1971--2018,0.44) (1993--2018,0.55) (2006--2018,0.62)};
\addlegendentry{الأنهار الجليدية}
\addplot[fill=omThm!55, draw=omThm!80!black] coordinates {(1971--2018,0.25) (1993--2018,0.43) (2006--2018,0.63)};
\addlegendentry{غرينلاند}
\addplot[fill=omThm!25, draw=omThm!70!black] coordinates {(1971--2018,0.14) (1993--2018,0.25) (2006--2018,0.37)};
\addlegendentry{القارّة القطبية الجنوبية}
\addplot[fill=black!20, draw=black!50] coordinates {(1971--2018,0.15) (1993--2018,0.31) (2006--2018,0.60)};
\addlegendentry{المياه البرّية}
\node[fill=black, rotate=45, minimum size=5pt, inner sep=0pt] at (axis cs:1971--2018,2.33) {};
\node[fill=black, rotate=45, minimum size=5pt, inner sep=0pt] at (axis cs:1993--2018,3.25) {};
\node[fill=black, rotate=45, minimum size=5pt, inner sep=0pt] at (axis cs:2006--2018,3.69) {};
\end{axis}
\end{tikzpicture}
\par\smallskip{\small \omterm{def:b2:anthropogenic-climate-change:sea-level}{ميزانية مستوى سطح البحر}: الإسهامات المرصودة في ارتفاع المتوسط العالمي لمستوى سطح
البحر، على هيئة معدّلات على مدى ثلاث فترات تنتهي عام 2018، مكدّسةً، والمعدّل المرصود (المعيّنات).
والفجوة بين الكدسة والمعيّن هي الجزء الذي لا تُغلقه الميزانية. القيم: تقرير التقييم السادس للهيئة الحكومية الدولية المعنية بتغيّر المناخ،
الفريق العامل الأول، الجدول 9.5.}
\end{omfigure}

\section*{التمارين}\label{sec:b2:anthropogenic-climate-change:exercises}

معطيات التمارين، ما لم يُذكر غير ذلك: تأثير مضاعفة ثنائي أكسيد الكربون $F_{2\times} =
\qty{3.93}{W.m^{-2}}$ انطلاقًا من $C_0 = \qty{278}{ppm}$؛ وثابت ستيفان--بولتزمان $\sigma =
\qty{5.670e-8}{W.m^{-2}.K^{-4}}$؛ ونصف قطر الأرض \qty{6371}{km}؛ والسنة الواحدة
\qty{3.156e7}{s}؛ والكتلتان الموليتان للكربون ولثنائي أكسيد الكربون 12.011 و\qty{44.009}{g.mol^{-1}}؛
والتغذيات الراجعة \omterm{def:b2:anthropogenic-climate-change:sea-level}{وميزانية مستوى سطح البحر} في الفصل.
% ledger: ipcc:erf-2xco2, const:sigma, earth:Rmean, ciaaw:abridged-2024, ar6ch7:table710, ar6ch9:table95

\begin{exercise}[$\star$]\label{exo:b2:anthropogenic-climate-change:1}
تأثير مضاعفة ثنائي أكسيد الكربون \qty{3.75}{W.m^{-2}} بعد تكيّف الستراتوسفير؛
ثم تتكيّف التروبوسفير بالمقادير $-0.60$ (درجة الحرارة)، و+0.22 (بخار الماء)، و+0.45 (السحب)
و+0.11 (سطح اليابسة) \unit{W.m^{-2}}. (أ) أوجد \omterm{def:b2:anthropogenic-climate-change:erf}{التأثير الإشعاعي الفعّال}. (ب) بأيّ نسبة مئوية يزيد على
القيمة \qty{3.71}{W.m^{-2}} في التقييم السابق؟ (ج) أيّ التكيّفات سالب، ولماذا تُنقص تروبوسفير أدفأ
التأثير؟
\end{exercise}
% ledger: ar6ch7:2xco2

\begin{exercise}[$\star$]\label{exo:b2:anthropogenic-climate-change:2}
اكتسب النظام المناخي \qty{434.9}{ZJ} من 1971 إلى 2018 و\qty{152.4}{ZJ} من 2006 إلى 2018،
محسوبةً بين منتصفي السنتين الأولى والأخيرة. (أ) أوجد مساحة سطح الأرض.
(ب) حوّل كلّ كسب إلى اختلال متوسط بوحدة \unit{W.m^{-2}}. (ج) أخذ المحيط
\qty{91}{\percent}: كم زيتاجولًا على مدى 1971--2018؟
\end{exercise}
% ledger: ar6ch7:inventory, earth:Rmean

\begin{exercise}[$\star$]\label{exo:b2:anthropogenic-climate-change:3}
(أ) اجمع القيم المركزية للتغذيات الراجعة في الشكل، وقارن بالقيمة الصافية. (ب) أوجد
\omterm{def:b2:anthropogenic-climate-change:ecs}{حساسية المناخ عند التوازن} بالقيمة الصافية، والاحترار \omterm{def:b2:anthropogenic-climate-change:feedbacks}{باستجابة بلانك} وحدها. (ج) بأيّ عامل تضخّم
التغذيات الراجعة الأخرى احترار بلانك؟
\end{exercise}

\begin{exercise}[$\star$]\label{exo:b2:anthropogenic-climate-change:4}
(أ) حوّل استجابة TCRE قدرها \qty{0.45}{\celsius} لكلّ \qty{1000}{Gt} من ثنائي أكسيد الكربون إلى درجات لكلّ
\qty{1000}{PgC} من الكربون، وافعل الشيء نفسه لمداها المرجح، من 0.27 إلى 0.63. (ب) حوّل
الكمّية \qty{2390}{Gt} من ثنائي أكسيد الكربون المنبعثة من 1850 إلى 2019 إلى بيتاغرامات من الكربون. (ج) ما
الاحترار الذي تعطيه الاستجابة TCRE لهذه الانبعاثات؟
\end{exercise}
% ledger: ar6spm:tcre, ar6spm:budget

\begin{exercise}[$\star\star$]\label{exo:b2:anthropogenic-climate-change:5}
(أ) يبعث السطح \qty{398}{W.m^{-2}} ويحتفظ أثر الدفيئة بالقيمة \qty{159}{W.m^{-2}} منها.
أوجد $g$ ودرجة حرارة سطح جسم أسود يبعث \qty{398}{W.m^{-2}}. (ب) أوجد
\omterm{def:b2:anthropogenic-climate-change:feedbacks}{استجابة بلانك} عند \qty{288}{K}، وعند \qty{298}{K} مع القيمة نفسها للكسر $g$. (ج) هل كانت \omterm{def:b2:anthropogenic-climate-change:feedbacks}{استجابة بلانك}
وحدها ستجعل أرضًا أدفأ أكثر حساسيةً أم أقلّ؟
\end{exercise}
% ledger: ipcc:budget-surface, ar6ch7:planck

\begin{exercise}[$\star\star$]\label{exo:b2:anthropogenic-climate-change:6}
تعطي الأقمار الصناعية تغذية راجعة لبخار الماء قدرها $1.85 \pm 0.32$ وتعطي الدراسات تغذية راجعة لمعدّل الانخفاض الحراري قدرها
$-0.50 \pm 0.20$ \unit{W.m^{-2}.\celsius^{-1}} (انحراف معياري واحد). (أ) أوجد مجموعهما.
(ب) لو كانتا مستقلّتين، فما الانحراف المعياري للمجموع؟ (ج) تعطي النماذج
المجموعَ انحرافًا معياريًا قدره \num{0.10} فقط. ماذا يقول ذلك عن الخطأين؟
\end{exercise}
% ledger: ar6ch7:wv-lr

\begin{exercise}[$\star\star$]\label{exo:b2:anthropogenic-climate-change:7}
من 1850--1900 إلى 2006--2019، $\Delta N = \qty{0.59}{W.m^{-2}}$، و$\Delta T = \qty{1.03}{\celsius}$
و$\Delta F = \qty{2.20}{W.m^{-2}}$. (أ) أوجد $\lambda$ والحساسية الفعلية. (ب) يُتوقّع أن يصبح
معامل التغذية الراجعة أقلّ سلبيةً بمقدار $\lambda' = \qty{0.5}{W.m^{-2}.\celsius^{-1}}$
مع تطوّر نمط الاحترار. أوجد الحساسية $-F_{2\times}/(\lambda + \lambda')$.
(ج) أعِد (ب) مع $\lambda' = 0.1$.
\end{exercise}
% ledger: ar6ch7:energy-budget, ar6ch7:pattern, ar6ch7:pattern-ecs

\begin{exercise}[$\star\star$]\label{exo:b2:anthropogenic-climate-change:8}
(أ) انطلاقًا من ارتفاع قدره \qty{0.617}{m} لكلّ درجة من احترار المحيط كلّه، ومتوسط عمق
\qty{3682}{m}، أوجد معامل التمدّد الفعلي. (ب) اكتسب المحيط \qty{396}{ZJ} على مدى
1971--2018. بكفاءة تمدّد قدرها \qty{0.113}{m} لكلّ \qty{e24}{J}، أوجد الارتفاع
الحراري، وقارنه بالقيمة \qty{47.5}{mm} في الميزانية. (ج) طبقة سمكها \qty{700}{m} تحترّ بمقدار
\qty{0.2}{\celsius} مع $\alpha = \qty{2e-4}{\celsius^{-1}}$: بكم ترفع مستوى سطح البحر؟
\end{exercise}
% ledger: ar6ch9:expansion, charette:ocean-volume, ar6ch7:inventory, ar6ch9:table95

\begin{exercise}[$\star\star$]\label{exo:b2:anthropogenic-climate-change:9}
بصمة نموذج، على مدى أربعة عقود، هي $x = (0.10, 0.20, 0.30, 0.40)$ \unit{\celsius}؛
والأرصاد هي $y = (0.12, 0.18, 0.33, 0.41)$ \unit{\celsius}. (أ) أوجد معامل القياس بالمربّعات الصغرى
$\beta$ في $y = \beta x + \varepsilon$. (ب) خطؤه المعياري، المقدّر من محاكاة
غير خاضعة لتأثير، هو \num{0.08}. هل الاستجابة مكتشَفة؟ وهل هي متوافقة مع النموذج؟ (ج) ماذا
تعني قيمة $\beta$ قدرها 0.5 مع الخطأ نفسه؟
\end{exercise}

\begin{exercise}[$\star\star\star$]\label{exo:b2:anthropogenic-climate-change:10}
بصمتان، لغازات الدفيئة $g = (0.2, 0.4, 0.6, 0.8)$ وللهباء الجوي $a = (-0.1, -0.2,
-0.2, -0.3)$، كلتاهما بوحدة \unit{\celsius} على مدى أربعة عقود، تُلاءَمان مع الأرصاد $y = (0.1,
0.2, 0.4, 0.5)$. (أ) اكتب المعادلات الطبيعية للملاءمة بالمربّعات الصغرى $y = \beta_g g + \beta_a a
+ \varepsilon$ وحلّها. (ب) أوجد الاحترار المُسنَد إلى كلّ محرّك في العقد الأخير.
(ج) لماذا يزداد الشكّ في $\beta_g$ و$\beta_a$ حين تتشابه البصمتان؟
\end{exercise}

\begin{exercise}[$\star\star\star$]\label{exo:b2:anthropogenic-climate-change:11}
استعمل طريقة \omterm{def:b2:anthropogenic-climate-change:budget}{ميزانية الكربون المتبقّية} مع استجابة TCRE قدرها \qty{0.45}{\celsius} لكلّ
\qty{1000}{Gt}، والكمّية \qty{210}{Gt} المنبعثة من 2015 إلى 2019 و\qty{26}{Gt} لكلّ درجة من أجل
التغذيات الراجعة لنظام الأرض. (أ) أوجد الميزانية لحدّ \qty{2}{\celsius}، أي \qty{0.93}{\celsius} فوق
2010--2019، إذا أضافت المحرّكات الأخرى \qty{0.1}{\celsius}. (ب) ما الاحترار الإضافي من المحرّكات
الأخرى الذي يجعلها مساوية للقيمة المقيّمة \qty{1350}{Gt}؟ هل يقع ضمن المدى المقيَّم؟ (ج) بيّن
أن كلّ \qty{0.1}{\celsius} من الاحترار غير الناجم عن CO$_2$ يغيّر أيّ ميزانية بالمقدار نفسه.
\end{exercise}

\begin{exercise}[$\star\star\star$]\label{exo:b2:anthropogenic-climate-change:12}
للتغذية الراجعة الصافية مدى مرجح جدًا من $-1.81$ إلى \qty{-0.51}{W.m^{-2}.\celsius^{-1}}، حول
$-1.16$. (أ) أوجد حساسية التوازن عند الطرفين وعند المركز. (ب) بيّن أن مدى للمعامل $\lambda$
متناظرًا حول مركزه يعطي مدى للحساسية غير متناظر، وقل أيّ جانبيه أطول. (ج) لماذا
تجد التقييمات استبعادَ حساسية مرتفعة أصعب من استبعاد حساسية منخفضة؟
\end{exercise}
% ledger: ar6ch7:table710

\begin{problem}[{مسألة نهاية الأسبوع --- قراءة الميزانية}]\label{pb:b2:anthropogenic-climate-change:1}
تتتبّع هذه المسألة طاقة التقييم وكربونه من التأثير إلى الميزانية
المتبقّية لحدّ \qty{1.5}{\celsius}. استعمل معطيات التمارين، و: \omterm{def:b2:anthropogenic-climate-change:erf}{التأثير الإشعاعي الفعّال} 1750--2019 (\unit{W.m^{-2}})
لثنائي أكسيد الكربون 2.16، والميثان 0.54، وأكسيد النيتروز 0.21، والغازات المهلجنة 0.41، والأوزون 0.47،
وبخار الماء الستراتوسفيري 0.05، والتفاعل بين الهباء الجوي والإشعاع $-0.22$، والتفاعل بين الهباء الجوي والسحب $-0.84$، واستخدام الأراضي
$-0.20$، والكربون الأسود على الثلج 0.08، والذيول التكاثفية للطائرات 0.06؛ وثنائي أكسيد الكربون \qty{410}{ppm} عام 2019
و\qty{427.35}{ppm} عام 2025؛ والتسخين \qty{0.79}{W.m^{-2}} على مدى 2006--2018؛ وتغيّرات ميزانية الطاقة
$\Delta N = 0.59$، و$\Delta F = \qty{2.20}{W.m^{-2}}$، و$\Delta T = \qty{1.03}{\celsius}$؛ والاستجابة TCRE
\qty{0.45}{\celsius} لكلّ \qty{1000}{Gt} من ثنائي أكسيد الكربون (على نحو مرجح من 0.27 إلى 0.63)؛ و\qty{2390}{Gt}
منبعثة في 1850--2019؛ والاحترار البشري \qty{1.07}{\celsius} للفترة 2010--2019؛ و\qty{210}{Gt} منبعثة من
2015 إلى 2019؛ والاحترار الإضافي غير الناجم عن CO$_2$ \qty{0.1}{\celsius}، وهو الطرف الأدنى للمدى المقيَّم من 0.1
إلى \qty{0.2}{\celsius}؛ والتغذيات الراجعة لنظام الأرض \qty{26}{Gt} لكلّ درجة. قرّب إلى ثلاثة أرقام
معنوية.
% ledger: ipcc:erf-agents, ipcc:erf-2xco2, data:co2-mlo, ar6ch7:inventory, ar6ch7:energy-budget, ar6spm:tcre, ar6spm:budget, ar6ch5:budget-parts

\noindent\textbf{الجزء الأول --- التأثير.}
\begin{enumerate}
  \item اجمع تأثير غازات الدفيئة الأربعة جيّدة الامتزاج.
  \item اجمع تأثير الهباء الجوي.
  \item اجمع كلّ الحدود البشرية. قارن بالقيمة المقيّمة \qty{2.72}{W.m^{-2}}.
  \item تحقّق من تأثير ثنائي أكسيد الكربون عام 2019 بالقانون اللوغاريتمي.
  \item أوجد تأثير ثنائي أكسيد الكربون عام 2025 بالقانون نفسه.
  \item أيّ الحدود له أوسع شكّ، ولماذا يهيمن على شكّ المجموع؟
\end{enumerate}

\noindent\textbf{الجزء الثاني --- ميزانية الطاقة.}
\begin{enumerate}[resume]
  \item كم من الطاقة اكتسب النظام المناخي على مدى 2006--2018، بالزيتاجول؟
  \item كم منها ذهب إلى المحيط؟
  \item أوجد معامل التغذية الراجعة من تغيّرات ميزانية الطاقة.
  \item أوجد حساسية المناخ الفعلية عند التوازن.
  \item أوجد \omterm{def:b2:anthropogenic-climate-change:tcr}{الاستجابة المناخية العابرة} من المعطيات نفسها.
  \item أوجد \omterm{def:b2:anthropogenic-climate-change:ecs}{حساسية المناخ عند التوازن} من التغذية الراجعة الصافية في الشكل، وقارن بين القيمتين.
  \item أعطِ السبب الفيزيائي لكون التقدير التاريخي هو الأدنى.
\end{enumerate}

\noindent\textbf{الجزء الثالث --- الاحترار لكلّ طنّ.}
\begin{enumerate}[resume]
  \item حوّل الاستجابة TCRE إلى درجات لكلّ \qty{1000}{PgC} من الكربون.
  \item حوّل مداها المرجح.
  \item حوّل انبعاثات 1850--2019 إلى بيتاغرامات من الكربون.
  \item ما الاحترار الذي تعطيه الاستجابة TCRE لانبعاثات 1850--2019؟
  \item قارن بالاحترار البشري البالغ \qty{1.07}{\celsius}. ماذا تقول المقارنة عن
        الأثر الصافي للمحرّكات البشرية الأخرى؟
\end{enumerate}

\noindent\textbf{الجزء الرابع --- الميزانية المتبقّية لحدّ \qty{1.5}{\celsius}.}
\begin{enumerate}[resume]
  \item ما الاحترار المتبقّي فوق 2010--2019؟
  \item اطرح الاحترار غير الناجم عن CO$_2$ واقسم على الاستجابة TCRE: ما الميزانية المتبقّية ابتداءً من منتصف
        2010--2019؟
  \item اطرح انبعاثات 2015--2019.
  \item اطرح التغذيات الراجعة لنظام الأرض على الاحترار المتبقّي.
  \item بكم تتغيّر الميزانية إذا أضافت المحرّكات الأخرى \qty{0.2}{\celsius} بدلًا من ذلك؟
  \item خلاصة القول: ما \omterm{def:b2:anthropogenic-climate-change:budget}{ميزانية الكربون المتبقّية} لحدّ \qty{1.5}{\celsius}، ابتداءً من مطلع
        2020، محسوبةً من الاستجابة TCRE؟ قارن بالقيمة المقيّمة \qty{500}{Gt}.
\end{enumerate}
\end{problem}
